%!TEX root = ./main.tex
%
% This file is part of the i10 thesis template developed and used by the
% Media Computing Group at RWTH Aachen University.
% The current version of this template can be obtained at
% <http://www.media.informatik.rwth-aachen.de/karrer.html>.

\documentclass[11pt, a4paper, titlepage]{book}

% ! ACHTUNG !
% Nach dem ersten LaTeX durchlauf auf der Kommandozeile
% makeindex -s main.ist main.idx
% ausf�hren, sonst wird der Index nicht so sch�n formatiert.
% Leider macht TeXShop den makeindex Aufruf nur ohne Parameter.

% !ATTENTION!
% After running LaTeX for the first time using this template call
% makeindex -s main.ist main.idx
% to get the correct formatting for the index.

%Pakete und eigene Befehle - am besten durchlesen!
%Includes all neccessary packets and contains some useful commands - read this file!
%!TEX root = ./main.tex

% This file is part of the i10 thesis template developed and used by the
% Media Computing Group at RWTH Aachen University.
% The current version of this template can be obtained at
% <http://www.media.informatik.rwth-aachen.de/karrer.html>.



%-----------------------------------------------------------------------------------------------------------------------------------------
% Befehle
% commands
%-----------------------------------------------------------------------------------------------------------------------------------------

%----------------------------------------------------------------------------------
% \myBigFigure	[ LABEL_PREFIX (optional) ]
%				{ FILENAME (without extension) }
%				{ CAPTION TEXT }
%				{ SHORT VERSION OF CAPTION TEXT }
%
%Bild wird in kompletter Breite gesetzt, die Kurzversion der Bildunterschrift erscheint im Abbildungsverzeichnis
%picture using full width of the page, the short caption is what appears in the list of figures index

%----------------------------------------------------------------------------------
% \myFrameBigFigure	[ LABEL_PREFIX (optional) ]
%					{ FILENAME (without extension) }
%					{ CAPTION TEXT }
%					{ SHORT VERSION OF CAPTION TEXT }
%
%Bild wird in kompletter Breite gesetzt und eingerahmt, die Kurzversion der Bildunterschrift erscheint im Abbildungsverzeichnis
%picture with frame using the full width of the page, the short caption is what appears in the list of figures index

%----------------------------------------------------------------------------------
% \myHUGEFigure	[ LABEL_PREFIX (optional) ]
%				{ FILENAME (without extension) }
%				{ CAPTION TEXT }
%				{ SHORT VERSION OF CAPTION TEXT }
%
%Bild wird rotiert und quer in kompletter Breite gesetzt, die Kurzversion der Bildunterschrift erscheint im Abbildungsverzeichnis
%landscape picture using the full width of the rotated page, the short caption is what appears in the list of figures index

%----------------------------------------------------------------------------------
% \myFigure	[ LABEL_PREFIX (optional) ]
%			{ FILENAME (without extension) }
%			{ CAPTION TEXT }
%			{ SHORT VERSION OF CAPTION TEXT }
%
%Bild wird in der Breite der textspalte gesetzt, die Kurzversion der Bildunterschrift erscheint im Abbildungsverzeichnis
%picture using the width of the text column, the short caption is what appears in the list of figures index

%----------------------------------------------------------------------------------
% \myImgRef	[ LABEL_PREFIX (optional) ]
%			{ LABEL OF THE IMAGE }
%
%referenziert das angegebene Bild
%reference to an image

%----------------------------------------------------------------------------------
% \myBigTable	{ YOUR TABULAR DEFINITION }
%			{ CAPTION TEXT }
%			{ TABLE_LABLE }
%
%Tabelle wird in kompletter Breite gesetzt
%table using the full width of the page

%----------------------------------------------------------------------------------
% \myTable	{ YOUR TABULAR DEFINITION }
%			{ CAPTION TEXT }
%			{ TABLE_LABLE }
%
%Tabelle wird in der Breite der Textspalte gesetzt
%table using the width of the text column

%----------------------------------------------------------------------------------
% \myTxtRef	{ LABLE }
%
%Referenz auf Kapitel oder Abschnitte - gibt nummer und namen aus, z.B.: 5.3---"Yaddahyaddah"
%references chapters or sections, outputs number and title, e.g., 5.3---"Yaddahyaddah"

%----------------------------------------------------------------------------------
% \myUnderscore
%
%Setzt einen "sch�nen" Unterstrich f�r URLs
%typesets a 'nice' underscore for URLs

%----------------------------------------------------------------------------------
%\myTilde
%
%Setzt eine "sch�ne" Tilde f�r URLs
%typesets a 'nice' tilde for URLs

%----------------------------------------------------------------------------------
% \myURL	{ TYPESET VERSION OF ANCHOR }
%			{ PRISTINE URL }
%			{ TYPESET VERSION OF URL }
%
%Setzt eine URL
%die typographisch "sch�ne" version erscheint in einer Fu�note,
%im Text erscheint der Ankertext, verlinkt ist die "echte" URL
%typesets a URL
%the typographically correct version appears as a footnote,
%the anchor appears in the text, the link points to the pristine URL

%----------------------------------------------------------------------------------
% \mySimpleURL	{ TYPESET VERSION OF ANCHOR }
%				{ PRISTINE URL }
%
%Setzt eine URL
%die URL erscheint in einer Fu�note,
%im Text erscheint der Ankertext, die URL ist verlinkt
%typesets a URL
%the URL appears as a footnote,
%the anchor appears in the text, the link points to the URL

%----------------------------------------------------------------------------------
% \myProjectURL	{ TYPESET VERSION OF ANCHOR }
%				{ PRISTINE URL INSIDE PROJECT DIRECTORY }
%				{ TYPESET VERSION OF URL INSIDE PROJECT DIRECTORY }
%
%Setzt eine URL innerhalb des Projektverzeichnisses auf "media"
%ACCOUNT muss durch den eigenen Usernamen ersetzt werden
%die typographisch "sch�ne" version erscheint in einer Fu�note,
%im Text erscheint der Ankertext, verlinkt ist die "echte" URL
%typesets a URL inside the project directory on 'media'
%replace ACCOUNT with your username
%the typographically correct version appears as a footnote,
%the anchor appears in the text, the link points to the pristine URL

%----------------------------------------------------------------------------------
% \mnote	{ MARGIN NOTE }
%
%Setzt eine Randnotitz
%puts a comment into the margin in small sans-serif font

%----------------------------------------------------------------------------------
% \todo	{ TODO MARGIN NOTE }
%
%Setzt eine "ToDo"-Randnotitz in rot zur Erinnerung
%puts a 'todo' comment into the margin in red

%----------------------------------------------------------------------------------
% \chapterquote	{ QUOTATION }
%				{ SOURCE }
%
%Setzt ein Zitat zum Einleiten eines Kapitels
%outputs a quote with its source, can be used as an introduction to chapters

%----------------------------------------------------------------------------------
% \myDefBox	{ TERM }
%			{ DEFINITION }
%
%Setzt eine Randnotitz und eine farbige Box (Textspaltenbreite),
%welche einen Begriff und seine Definition enth�lt
%outputs a margin note and a colored box (width of the text column) containing a term and its definition

%----------------------------------------------------------------------------------
% \myBigDefBox	{ TERM }
%				{ DEFINITION }
%
%Setzt eine farbige Box (Seitenbreite), welche einen Begriff und seine Definition enth�lt
%outputs a colored box (width of the page) containing a term and its definition

%----------------------------------------------------------------------------------
% \myDownloadURL	{ TYPESET DOWNLOAD NAME }
%					{ PRISTINE VERSION OF FILENAME }
%					{ TYPESET VERSION OF FILENAME }
%
%Setzt eine farbige Box, welche einen Downloadlink enth�lt
%outputs a colored box containing a download link

%----------------------------------------------------------------------------------
% \emptydoublepage
%
% Leere Doppelseite ohne Kopf- oder Fu�zeile am Ende von Kapiteln
% Clear double page without any header or footer at end of chapters

%----------------------------------------------------------------------------------
% \pagebreak	[ SOME STRANGE LATEX VALUE ]
%
%Eklige pagebreaks f�r den Druck (falls es nicht mehr anders geht)
%pagebreaks for the final print version (last resort weapon against wrong pagebreaks by LaTeX)

%----------------------------------------------------------------------------------
% \TM
%
%Setzt ein (TM) Symbol
%Places a (TM) symbol


%----------------------------------------------------------------------------------
%Packages and parameters
%----------------------------------------------------------------------------------

%Inputencoding f�r den Mac
%inputencoding for the mac
\usepackage[utf8]{inputenc}

%Deutschpaket fŸr das Proseminar
\usepackage{german}
%Mathe- und Symbolpakete

%packages for mathematical symbols
\usepackage{latexsym}
\usepackage{amsmath}
\usepackage{amssymb}

%Tabellengestaltung
%table design
\usepackage{booktabs}

%Grafikpaket
%grahics package
\usepackage{color,graphicx}

%relativer Pfad zu den Bildern
%path to your image folder
\graphicspath{{images/}}

%Abs�tze werden nicht eingezogen, sondern vertikal abgesetzt
%do not indent at new paragraphs but add a vertical offset
\usepackage{noindent}

%Palatino+Helvetica statt Computer Modern als standard fonts:
%change standard fonts to Palatino and Helvetica
\usepackage{palatino}

%Bibliographieeinstellungen
%bibliography settings
\usepackage{natbib}
\bibliographystyle{plainnat}

%Zitierbefehle
%citation commands
\newcommand{\fullcite}{\citep} %for "Author [1980]"
\renewcommand{\citeyear}{\citeyearpar} %for "[1980]"

%paket f�r erweiterte kontrollstrukturen
%package for control structures
\usepackage{ifthen}

%marginpar hack --- alle Randnotitzen sollten dann auf der richtigen Seite stehen
%marginpar hack --- moves margin notes to correct position
\usepackage{mparhack}

%lesbare verweise
%make readable references
\usepackage[pdftex,plainpages=false,pdfpagelabels]{hyperref}


%---------------------<Layout in the style of "A Pattern Approach to Interaction Design>---------------------------

% Change page headers and footers:
\usepackage{fancyhdr}
\pagestyle{fancy}
\fancyhf{}
\fancyhead[RE]{\slshape \nouppercase{\leftmark}}    % Even page header: "page   chapter"
\fancyhead[LO]{\slshape \nouppercase{\rightmark}}   % Odd  page header: "section   page"
\fancyhead[RO,LE]{\bfseries \thepage} 
\renewcommand{\headrulewidth}{1pt}    % Underline headers
\renewcommand{\footrulewidth}{0pt}    

\fancypagestyle{plain}{               % No chapter+section on chapter start pages
\fancyhf{}
\fancyhead[RO,LE]{\bfseries \thepage}
\renewcommand{\headrulewidth}{1pt}
\renewcommand{\footrulewidth}{0pt}
}

% Left headings: "1  INTRODUCTION"
\renewcommand{\chaptermark}[1]{%
\markboth{\thechapter\ \ \ \ #1}{}}

% Right headings: "1.1  Basics"
\renewcommand{\sectionmark}[1]{%
\markright{\thesection\ \ \ \ #1}{}}

% some Fancyhdr problem...
\addtolength{\headheight}{2pt} % To avoid overfull vboxes from fancyhdr


%creating a better way to change the layout for the abstract pages
\usepackage{geometry}

\ifthenelse{\lengthtest{\paperheight=250mm}}%
{% -----------------B5 Layout-----------------
% Page layout

%\pdfpageheight250mm
%\pdfpagewidth176mm
\geometry{	b5paper,
			top = 27mm,
			footskip = 10mm,
			inner = 19mm,
			outer = 39mm,
			textheight = 175mm,
			textwidth = 84mm,
			marginparsep = 3mm,
			marginparwidth = 32mm
}
\savegeometry{myText}
% -----------------/ B5 Layout-----------------
}%
{% -----------------A4 Layout-----------------
% Page layout
\geometry{	a4paper,
			twoside,
			includemp,
			includehead,
			top = 30mm,
			headsep = 10mm,
			bindingoffset = 10mm,
			inner = 20mm,
			outer = 40mm,
			bottom = 45mm,
			marginparsep = 10mm,
			marginparwidth = 30mm
}
\savegeometry{myText}
% -----------------/ A4 Layout-----------------
}
% Abstract layout
\geometry{	marginparsep = 0mm,
			marginparwidth = 0mm
}
\savegeometry{myAbstract}
\loadgeometry{myText}

\newlength{\fullwidth} % Width of text plus margin notes
\setlength{\fullwidth}{\textwidth}
\addtolength{\fullwidth}{\marginparsep}
\addtolength{\fullwidth}{\marginparwidth}

\setlength{\headwidth}{\fullwidth} % Header stretches over margin notes


%---------------------</Layout in the style of "A Pattern Approach to Interaction Design>---------------------------


%wird f�r die fl�chendeckende Ausgabe der Titelseite ben�tigt
%needed for the full-face titlepage
\usepackage{eso-pic}

%index verwenden
%make an index
\usepackage{makeidx}
\makeindex

%Index Formatierungshilfen
%formatting helpers for the index
\newcommand{\uu}[1]{\underline{#1}}
\newcommand{\ii}[1]{\textit{#1}}

%neue Definition der Index Umgebung
%redesign of the index
\renewenvironment{theindex}{%
  \vspace*{50pt}%
  {\Huge\bfseries\indexname}\par%
  \vspace*{40pt}%
  \setlength{\parskip}{0pt}%
  \setlength{\parindent}{0pt}%
  \small%
  \renewcommand{\item}{\par{}}%
  \renewcommand{\subitem}{\par\hspace{2em}- }%
}%
{}

%Maximale Gliederungstiefe, die noch ins Inhaltsverzeichnis aufgenommen wird
%maximum depth for the table of contents
\setcounter{tocdepth}{3}

%Vorschlag f�r ein sch�nes Farbschema
%Set of colors which look nice together
\usepackage{color}
\definecolor{orange_light}{rgb}{1,0.8,0.4}
\definecolor{orange_med}{rgb}{0.753,0.62,0.373}
\definecolor{orange_dark}{rgb}{0.506,0.412,0.251}

\definecolor{green_light}{rgb}{0.8,1,0.4}
\definecolor{green_med}{rgb}{0.635,0.745,0.376}
\definecolor{green_dark}{rgb}{0.435,0.498,0.255}

\definecolor{blue_light}{rgb}{0.4,0.8,1}
\definecolor{blue_med}{rgb}{0.365,0.624,0.749}
\definecolor{blue_dark}{rgb}{0.251,0.42,0.502}

\definecolor{pink_light}{rgb}{1,0.435,0.812}
\definecolor{pink_med}{rgb}{0.745,0.38,0.62}
\definecolor{pink_dark}{rgb}{0.498,0.255,0.416}

\definecolor{yellow_light}{rgb}{1,1,0.4}
\definecolor{yellow_med}{rgb}{0.757,0.745,0.373}
\definecolor{yellow_dark}{rgb}{0.506,0.49,0.251}

%blau (f�r URLs)
%blue (for URLs)
\definecolor{blue}{rgb}{0,0,1}

%notwendig f�r die korrekte Erkennung, auf welcher Seite sich eine Abbildung befindet.
%we need this to determine if a figure is on an odd or even page
\usepackage{chngpage}

%Hiermit k�nnen die Abbildungslegenden frei gestaltet werden
%we need this to redesign the captions
\usepackage[font=normalsize,labelfont=bf]{caption}

%Abbildungen kommen auf eine eigene Seite, wenn sie mehr als 85% des Platzes
%auf einer Seite einnehmen
%if a figure takes more than 85% of a page it will be typeset on a separate page
\renewcommand{\floatpagefraction}{0.85}

%Ben�tigt um gro�e Abbildungen gedreht auf eine Seite zu setzen
%we need this to rotate big figures
\usepackage[figuresright]{rotating}

%Verschiedene L�ngenma�e f�r Textboxen
%dimensions for textboxes
\newlength{\myDefBoxWidth}
\setlength{\myDefBoxWidth}{\textwidth}
\addtolength{\myDefBoxWidth}{-4mm}
\newlength{\myBigDefBoxWidth}
\setlength{\myBigDefBoxWidth}{\fullwidth}
\addtolength{\myBigDefBoxWidth}{-4mm}

%Formathilfen f�r MatLab Code (wer's braucht...)
%pre-defined matlab code formats
\usepackage{alltt}
\definecolor{string}{rgb}{0.7,0.0,0.0}
\definecolor{comment}{rgb}{0.13,0.54,0.13}
\definecolor{keyword}{rgb}{0.0,0.0,1.0}



%-----------------------------------------------------------------------------------------------------------------------------------------
% neue Befehle
% new commands
%-----------------------------------------------------------------------------------------------------------------------------------------

%----------------------------------------------------------------------------------
% \myBigFigure	[ LABEL_PREFIX (optional) ]
%				{ FILENAME (without extension) }
%				{ CAPTION TEXT }
%				{ SHORT VERSION OF CAPTION TEXT }
%
%Bild wird in kompletter Breite gesetzt
%picture using full width of the page
\newcommand{\myBigFigure}[4][image]
{%
\begin{figure}[t!bp]%
	\checkoddpage%
	\ifcpoddpage%
		%nothing
	\else
		\hspace{-\marginparsep}\hspace{-\marginparwidth}%
	\fi
	%use minipage to center the label beneath the figure
	\begin{minipage}{\fullwidth}%
		\includegraphics[width= \fullwidth]{#2}%
		\caption[#4]{#3}%
		\label{#1_#2}%
	\end{minipage}%
\end{figure}%
}


%----------------------------------------------------------------------------------
% \myFrameBigFigure	[ LABEL_PREFIX (optional) ]
%					{ FILENAME (without extension) }
%					{ CAPTION TEXT }
%					{ SHORT VERSION OF CAPTION TEXT }
%
%Bild wird in kompletter Breite gesetzt und eingerahmt
%picture with frame using the full width of the page
\newcommand{\myFrameBigFigure}[4][image]
{
\begin{figure}[t!bp]
	\checkoddpage
	\ifcpoddpage
		%nothing
	\else
		\hspace{-\marginparsep}\hspace{-\marginparwidth}
	\fi
	%use minipage to center the label beneath the figure
	\begin{minipage}{\fullwidth}
	\frame{%
		\includegraphics[width= \fullwidth]{#2}%
		}
		\caption[#4]{#3}
		\label{#1_#2}
	\end{minipage}
\end{figure}
}

%----------------------------------------------------------------------------------
% \myHUGEFigure	[ LABEL_PREFIX (optional) ]
%				{ FILENAME (without extension) }
%				{ CAPTION TEXT }
%				{ SHORT VERSION OF CAPTION TEXT }
%
%Bild wird rotiert und quer in kompletter Breite gesetzt
%landscape picture using the full width of the rotated page
\newcommand{\myHugeFigure}[4][image]
{
\begin{sidewaysfigure}[t!bp]
	\checkoddpage
	\ifcpoddpage
		%nothing
		\vspace{\marginparsep}\vspace{\marginparwidth}
	\else
		%nothing
		\vspace{-\marginparsep}\vspace{-\marginparwidth}
	\fi
		\includegraphics[width= \textheight]{#2}
		\caption[#4]{#3}
		\label{#1_#2}
	
\end{sidewaysfigure}
}

%----------------------------------------------------------------------------------
% \myFigure	[ LABEL_PREFIX (optional) ]
%			{ FILENAME (without extension) }
%			{ CAPTION TEXT }
%			{ SHORT VERSION OF CAPTION TEXT }
%
%Bild wird in der Breite der textspalte gesetzt
%picture using the width of the text column
\newcommand{\myFigure}[4][image]%
{%
\begin{figure}[t!bp]%
	\begin{center}%
		\includegraphics[width= \textwidth]{#2}%
		\caption[#4]{#3}
		\label{#1_#2}%
	\end{center}%
\end{figure}%
}%

%----------------------------------------------------------------------------------
% \myImgRef	[ LABEL_PREFIX (optional) ]
%			{ LABEL OF THE IMAGE }
%
%referenziert das angegebene Bild
%reference to an image
\newcommand{\myImgRef}[2][image]%
{%
	\ref{#1_#2}%
}%

%----------------------------------------------------------------------------------
% \myBigTable	{ YOUR TABULAR DEFINITION }
%			{ CAPTION TEXT }
%			{ TABLE_LABLE }
%
%Tabelle wird in kompletter Breite gesetzt
%table using the full width of the page
\newcommand{\myBigTable}[3]%
{%
\begin{table}[htdp]%
	\checkoddpage%
	\ifcpoddpage%
		%nothing
	\else%
		\hspace{-\marginparsep}\hspace{-\marginparwidth}%
	\fi%
	\begin{minipage}{\fullwidth}%
		\begin{center}%
			#1%
			\caption{#2}%
			\label{#3}%
		\end{center}%	
	\end{minipage}%
\end{table}%
}%

%----------------------------------------------------------------------------------
% \myTable	{ YOUR TABULAR DEFINITION }
%			{ CAPTION TEXT }
%			{ TABLE_LABLE }
%
%Tabelle wird in der Breite der Textspalte gesetzt
%table using the width of the text column
\newcommand{\myTable}[3]%
{%
\begin{table}[htdp]%
	\begin{center}%
		#1%
		\caption{#2}%
		\label{#3}%
	\end{center}%	
\end{table}%
}%

%----------------------------------------------------------------------------------
% \myTxtRef	{ LABLE }
%
%Referenz auf Kapitel oder Abschnitte - gibt nummer und namen aus, z.B.: 5.3---"Yaddahyaddah"
%references chapters or sections, outputs number and title, e.g., 5.3---"Yaddahyaddah"
\newcommand{\myTxtRef}[1]
{%
	\ref{#1} ``\nameref{#1}''%
}

%----------------------------------------------------------------------------------
% \myTxtRefPP	{ LABLE }
%
%Referenz auf Kapitel oder Abschnitte - gibt nummer, namen und seiten aus, z.B.: 5.3---"Yaddahyaddah" (p. 45)
%references chapters or sections, outputs number and title, e.g., 5.3---"Yaddahyaddah"
\newcommand{\myTxtRefPP}[1]
{%
	\ref{#1} ``\nameref{#1}'' (p.~\pageref{#1})%
}

%----------------------------------------------------------------------------------
% \myUnderscore
%
%Setzt einen "sch�nen" Unterstrich f�r URLs
%typesets a 'nice' underscore for URLs
\newcommand{\myUnderscore}{$\underline{\hspace{0.5em}}$}

%----------------------------------------------------------------------------------
%\myTilde
%
%Setzt eine "sch�ne" Tilde f�r URLs
%typesets a 'nice' tilde for URLs
\newcommand{\myTilde}{$\sim$}

%----------------------------------------------------------------------------------
% \myURL	{ TYPESET VERSION OF ANCHOR }
%			{ PRISTINE URL }
%			{ TYPESET VERSION OF URL }
%
%Setzt eine URL
%die typographisch "sch�ne" version erscheint in einer Fu�note,
%im Text erscheint der Ankertext, verlinkt ist die "echte" URL
%typesets a URL
%the typographically correct version appears as a footnote,
%the anchor appears in the text, the link points to the pristine URL
\newcommand{\myURL}[3]%
{%
	\textcolor{blue}{%
		\href{#2}{#1}%
	}%
	\footnote{#3}%
}

%----------------------------------------------------------------------------------
% \mySimpleURL	{ TYPESET VERSION OF ANCHOR }
%				{ PRISTINE URL }
%
%Setzt eine URL
%die URL erscheint in einer Fu�note,
%im Text erscheint der Ankertext, die URL ist verlinkt
%typesets a URL
%the URL appears as a footnote,
%the anchor appears in the text, the link points to the URL
\newcommand{\mySimpleURL}[2]%
{%
	\textcolor{blue}{%
		\href{#2}{#1}%
	}%
	\footnote{#2}%
}

%----------------------------------------------------------------------------------
% \myProjectURL	{ TYPESET VERSION OF ANCHOR }
%				{ PRISTINE URL INSIDE PROJECT DIRECTORY }
%				{ TYPESET VERSION OF URL INSIDE PROJECT DIRECTORY }
%
%Setzt eine URL auf hci/public wo die Inhalte des WebServer Ordners auf "oliver" verf�gbar sind
%die typographisch "sch�ne" version erscheint in einer Fu�note,
%im Text erscheint der Ankertext, verlinkt ist die "echte" URL
%typesets a URL to hci/public from where the contents of the WebServer folder from oliver can be accessed
%the typographically correct version appears as a footnote,
%the anchor appears in the text, the link points to the pristine URL
\newcommand{\myProjectURL}[3]%
{%
	\textcolor{blue}{%
		\href{http://hci.rwth-aachen.de/public/#2}{#1}%
	}%
	\footnote{http://hci.rwth-aachen.de/public/#3}%
}

%----------------------------------------------------------------------------------
% \mnote	{ MARGIN NOTE }
%
%Setzt eine Randnotitz
%puts a comment into the margin in small sans-serif font
\newcommand{\mnote}[1]%
{%
	\leavevmode%
	\checkoddpage%
	\ifcpoddpage%
		\marginpar{\raggedright\textsf{{\footnotesize{#1}}}}%
	\else%
		\marginpar{\raggedleft\textsf{{\footnotesize{#1}}}}%
	\fi%
}
	
% leavevmode allows mnotes to be aligned with the first line of a paragraph
% NOTE: you have to put a "%" at the end of the line with the mnote, or you will get an extra blank at the beginning of the paragraph!

%----------------------------------------------------------------------------------
% \todo	{ TODO MARGIN NOTE }
%
%Setzt eine "ToDo"-Randnotitz in rot zur Erinnerung
%puts a 'todo' comment into the margin in red
\definecolor{red}{rgb}{1,0,0}
\newcommand{\todo}[1]{\mnote{\textcolor{red}{ToDo: #1}}}

%----------------------------------------------------------------------------------
% \chapterquote	{ QUOTATION }
%				{ SOURCE }
%
%Setzt ein Zitat zum Einleiten eines Kapitels
%outputs a quote with its source, can be used as an introduction to chapters
\newcommand{\chapterquote}[2]{
\begin{quotation}
    \begin{flushright}
	\noindent\emph{``{#1}''\\[1.5ex]---{#2}}
    \end{flushright}
\end{quotation}
}

%----------------------------------------------------------------------------------
% \myDefBox	{ TERM }
%			{ DEFINITION }
%
%Setzt eine Randnotitz und eine farbige Box (Textspaltenbreite),
%welche einen Begriff und seine Definition enth�lt
%outputs a margin note and a colored box (width of the text column) containing a term and its definition
\newcommand{\myDefBox}[2]
{%
	\setlength{\fboxrule}{1mm}%
	\fcolorbox{orange_med}{orange_light}%
	{%
		\parbox{\myDefBoxWidth}{{\bfseries\scshape#1:}\\#2}%
	}%
	\mnote{Definition:\\\emph{#1}}
}

%----------------------------------------------------------------------------------
% \myBigDefBox	{ TERM }
%				{ DEFINITION }
%
%Setzt eine farbige Box (Seitenbreite), welche einen Begriff und seine Definition enth�lt
%outputs a colored box (width of the page) containing a term and its definition
\newcommand{\myBigDefBox}[2]
{%
	\begin{figure}[h!]
	\setlength{\fboxrule}{1mm}%
	\checkoddpage%
	\ifcpoddpage%
		%nothing
	\else%
		\hspace{-\marginparsep}\hspace{-\marginparwidth}%
	\fi%
	\fcolorbox{orange_med}{orange_light}%
	{%
		\parbox{\myBigDefBoxWidth}{{\bfseries\scshape#1:}\\#2}%
	}%
	\end{figure}
}

%----------------------------------------------------------------------------------
% \myDownloadURL	{ TYPESET DOWNLOAD NAME }
%					{ PRISTINE VERSION OF FILENAME }
%					{ TYPESET VERSION OF FILENAME }
%
%Setzt eine farbige Box, welche einen Downloadlink enth�lt
%outputs a colored box containing a download link
\newcommand{\myDownloadURL}[3]{%
\checkoddpage%
	\ifcpoddpage%
		%nothing
	\else%
		\hspace{-\marginparsep}\hspace{-\marginparwidth}%
	\fi%
\setlength{\fboxrule}{1mm}%
\fcolorbox{green_med}{green_light}{%
\begin{minipage}{\myBigDefBoxWidth}%
\begin{center}%
\myProjectURL{#1}{folder/#2}{folder/#3}%
\end{center}%
\end{minipage}%
}%
}

%----------------------------------------------------------------------------------
% \emptydoublepage
%
% Leere Doppelseite ohne Kopf- oder Fu�zeile am Ende von Kapiteln
% Clear double page without any header or footer at end of chapters
\newcommand{\emptydoublepage}{\clearpage\thispagestyle{empty}\cleardoublepage}

%----------------------------------------------------------------------------------
% \pagebreak	[ SOME STRANGE LATEX VALUE ]
%
%Eklige pagebreaks f�r den Druck (falls es nicht mehr anders geht)
%pagebreaks for the final print version (last resort weapon against wrong pagebreaks by LaTeX)
\newcommand{\PB}[1][3]
{%
	\pagebreak[#1]%
}



%----------------------------------------------------------------------------------
% \TM
%
%Setzt ein (TM) Symbol
%Places a (TM) symbol
\newcommand{\TM}
{%
	\textsuperscript{\texttrademark}%
}



%--------------------------------------------------------------
%Dokumentspezifisches
%Stuff regarding your specific document
%--------------------------------------------------------------

%Trennungshilfen
%Hyphenation patterns
\hyphenation{
dieseswortwirdnichtgetrennt
diesesauchnicht
thiswordwillstayinoneline
thistoo
}


%--------------------------------------------------------------
\begin{document}

% gr��ere Wortabst�nde zulassen, um Trennungen zu vermeiden
% allow more flexible whitespaces to avoid hyphenation and overfull hboxes
\sloppy

% "see" Eintr�ge f�r den Index
% 'see'-entries for the index
\include{indexAddOn}


%--------------------------------------------------------------
\frontmatter

%Die Titelseite aus dem "images"-Verzeichnis wird verwendet
%use the titlepage from the 'images' directory 
\begin{titlepage}
\AddToShipoutPicture*{
\put(0,0){
\includegraphics[width=\paperwidth]{titlepage}}}
\strut
\end{titlepage}

\thispagestyle{empty}
\emptydoublepage

\tableofcontents
\emptydoublepage
%
\listoffigures
\emptydoublepage
%
\listoftables
\emptydoublepage

%Die �bersicht sollte in englischer und in deutscher Sprache verfasst werden (beim Proseminar nur deutsch)
%the abstract should include an english and a german version (only German in case of a Proseminar paper)
%# -*- coding: utf-8-unix -*-
%%==================================================

\begin{abstract}
本项目为年产50万吨MTO工厂的初步设计。通过分析当前国内外MTO生产和研究现状，对生产工艺进行了选择论证。然后运用Aspen软件模拟初步的工艺流程，并通过对一系列工艺参数，如精馏塔的塔板数—产品纯度、进料塔板数—产品纯度、产品纯度—回流比、再沸器负荷—回流比等进行灵敏度分析，优化设备操作条件，提高工艺的合理性和经济性。本设计还针对工艺流程进行换热网络设计和对全局换热网络进行了优化和评估，通过内部流股之间相互换热以减少公用工程的消耗，最终优化后节约$79.4\%$的热公用工程资源和$73.7\%$的冷公用工程资源。本设计还运用水夹点技术优化了用水网络，根据水硬度分类处理水操作单元，并合理再生利用，使得本项目新鲜水用量和废水排放量达到最小，优化后的用水网络节约用水$53.59\%$。本设计对于MTO工厂的生产和设计建造具有一定的现实指导意义。\\

\keywords{\zihao{-4} 工厂\quad 设计\quad MTO \quad 工艺 \quad 水夹点  \quad 网络 \quad 控制}
\end{abstract}

\begin{englishabstract}

This project is the preliminary design of a MTO plant with an annual output of 500,000 tons of light olefins. Based on the current production and research situation all through the world, the production method was selected and demonstrated. Aspen software was used to simulate the preliminary process. Heat integration method was applied to optimize the heat exchange network. Rational heat exchange between process streams were suggested which resulted in the decreasing of utilities consumption and exchanger number. The heat integration leaded to energy saving of $79.4\%$ of heat utilities and $73.7\%$ of the cold utilities. In addition, the water pinch technology was also implemented to optimize the water network. The water operating unit was classified according to water hardness, with a reasonable recycling. The amount of fresh water consumption and wastewater emission was minimized. The optimized water network achieved $53.59\%$ water saving. Finally, a preliminary economic analysis to the entire project was estimated in order to get the project construction cost and profitability. In summary, this design is of some practical significance for the production and design of the MTO industry.

\englishkeywords{\zihao{-4} Plant design\;Sensitivity analysis  \; Energy balance\; calculation \; Water pinch  Dynamic control}
\end{englishabstract}


\emptydoublepage

%In der Arbeit verwendete Konventionen (Formelsatz, Definitionen, etc.)
%conventions applied in the thesis (AE/BE, definitions, etc.)
%!TEX root = ./main.tex
%
% This file is part of the i10 thesis template developed and used by the
% Media Computing Group at RWTH Aachen University.
% The current version of this template can be obtained at
% <http://www.media.informatik.rwth-aachen.de/karrer.html>.

\chapter*{Konventionen\markboth{Konventionen}{Konventionen}}
\addcontentsline{toc}{chapter}{\protect\numberline{}Konventionen}

In dieser Ausarbeitung gelten folgende Konventionen.



\bigskip

\emph{Textkonventionen}

Definitionen von Fachbegriffen oder kurze Exkurse werden in farbige Rahmen gesetzt.

\myDefBox{Exkurs}
{Ein Exkurs ist eine in den Text eingef\"ugte oder als Anhang beigef\"ugte selbstst\"andige und in sich geschlossene Abschweifung.}

\medskip

Quellcode und Implementierungssyntax ist in Schreibmaschinenschrift gesetzt.

\texttt{myClass}

\medskip

Download links sind in farbigen Rahmen gesetzt.

\myDownloadURL{File: myFile}{file_number.file}{file\myUnderscore number.file}

%\pagebreak

\emph{Mathematische Konventionen}

Was auch immer...




\emptydoublepage

%--------------------------------------------------------------
\mainmatter


\lettrine[lines=3]{T}{he MC methodology} mainly considers two types of point-based temporal logics (PTLs) as the property specification language---\emph{linear} and \emph{branching}---which differ 
in the underlying model of time.
In linear PTLs, such as $\LTL$~\cite{Pnu77}, each moment in time has a unique possible future:
formulas are interpreted over (infinite) \emph{paths} of a Kripke structure, and thus they refer to a single computation of a system.
In branching PTLs,
such as $\CTL$ and $\CTLStar$~\cite{EH86}, each moment in time may evolve into several possible futures:
formulas are interpreted over \emph{states} of the Kripke structure, hence referring to all the possible system computations.

In the previous chapters we have assumed a \emph{state-based} semantics for $\HS$, which
induces a branching reference both in the future and in the past:
intervals/traces 
are \lq\lq forgetful\rq\rq{} of the history leading to their initial state, and  
the initial (resp., final) state of an interval may feature several predecessors (resp., successors).
A graphical account of the state-based semantics can be found in Figure~\ref{fig:ST}; a detailed explanation will be given in the following.

\begin{figure}[tp]
\centering
\begin{minipage}{0.36\linewidth}
\begin{tikzpicture}
				%				[->,>=stealth',shorten >=1pt,auto,node distance=0,main node/.style={circle,draw}]
				\filldraw [gray] (0,2) circle (2pt)
				(2,2) circle (2pt)
				(4,3) circle (2pt)
				(4,2) circle (2pt)
				(4,1) circle (2pt);
				%
					\filldraw [gray] (4,-1) circle (2pt)
				(2,-1) circle (2pt)
				(0,0) circle (2pt)
				(0,-1) circle (2pt)
				(0,-2) circle (2pt);
				%
				\draw [black]  (0,2) -- (4,2);
				\draw [black]  (0,-1) -- (4,-1);
				%
				\draw [black] (2,2) -- (4,3);
				\draw [black] (2,2) -- (4,1);
				\draw [black] (0,0) -- (2,-1);
				\draw [black] (0,-2) -- (2,-1);
				%
				\draw [dashed, orange] (0,2.2) -> (2,2.2) -> (4,3.2);
				\draw [dashed, red] (0,1.8) -> (2,1.8);
				\draw [dashed, orange] (0,0.2) -> (2,-0.8) -> (4,-0.8);
				\draw [dashed, red] (2,-1.2) -> (4,-1.2);
				
	%		{\tiny
				\node [orange] at (0.5,2.5) {$\varphi_1$};	
				\node [red] at (0.5,1.5) {$\hsBt\varphi_1$};
			\node [orange] at (2.5,-0.5) {$\varphi_1$};	
				\node [red] at (2.5,-1.5) {$\hsEt\varphi_1$};
		%			\node (b0) at (4,-0.5) {$\rho'=\rho(0,i)\cdot \rho^{(j+1)]}$};	
		%			\node (a1) at (1.5,0.2) {$\rho(i)=$};
		%			\node (a2) at (2,0.2) {$\rho(j)$};
		%			\node (a3) at (1.65,0.5) {$Pattern(\rho,i)= Pattern(\rho,j)$};
		%			\node (a4) at (4.6,0.5) {$Pattern(\rho,k) =\{ p \in \Prop : \Ku,\rho^{k]} \models p\}$};
						
	%			}
				
			\end{tikzpicture}
\end{minipage}	
\hfill 	
\begin{minipage}{0.54\linewidth}
			\begin{tikzpicture}
				%				[->,>=stealth',shorten >=1pt,auto,node distance=0,main node/.style={circle,draw}]
				%
					\filldraw [gray] (4,-1) circle (2pt)
				(2,-1) circle (2pt)
				(0,0) circle (2pt)
				(0,-1) circle (2pt)
				(0,-2) circle (2pt)
				(6,0) circle (2pt)
				(6,-1) circle (2pt)
				(6,-2) circle (2pt)	;
				%
				\draw [black]  (0,-1) -- (6,-1);
				%
				\draw [black] (4,-1) -- (6,0);
				\draw [black] (4,-1) -- (6,-2);
					\draw [black] (0,0) -- (2,-1);
				\draw [black] (0,-2) -- (2,-1);
				%
				\draw [dashed, orange] (0,0.2) -> (2,-0.8) -> (4,-0.8);
				\draw [dashed, orange] (2,-1.2) -> (4,-1.2) -> (6,-2.2);
				
	%		{\tiny
				\node [orange] at (1,0) {$\varphi_1$};	
				\node [red] at (2.5,-0.5) {$\hsAt\varphi_1$};
			\node [orange] at (5,-2) {$\varphi_2$};	
				\node [red] at (2.5,-1.5) {$\hsA \varphi_2$};
		%			\node (b0) at (4,-0.5) {$\rho'=\rho(0,i)\cdot \rho^{(j+1)]}$};	
		%			\node (a1) at (1.5,0.2) {$\rho(i)=$};
		%			\node (a2) at (2,0.2) {$\rho(j)$};
		%			\node (a3) at (1.65,0.5) {$Pattern(\rho,i)= Pattern(\rho,j)$};
		%			\node (a4) at (4.6,0.5) {$Pattern(\rho,k) =\{ p \in \Prop : \Ku,\rho^{k]} \models p\}$};
						
	%			}
				
			\end{tikzpicture}
\end{minipage}
%
    \caption{State-based semantic variant $\HS_\stat$: past and future are branching}
    \label{fig:ST}
\end{figure}

However, as we already said, $\HS$ MC has been simultaneously and independently studied also by Lomuscio and Michaliszyn in~\cite{LM13,LM14,lm16}:
there, the considered $\HS$ fragments are interpreted over the unwinding of a Kripke structure (\emph{computation-tree-based} semantics---see Figure~\ref{fig:CT}). %, and interval labeling takes into account only the endpoints of intervals. 
This induces a linear reference in the past (the initial state of an interval may feature only one predecessor), but branching in the future (the final state features several successors). Moreover, the computation history is never forgotten and increases with time.

\input{Chaps/TOCL17/Imgs/semPicturesCT.tex}

In this chapter, we study the expressiveness of $\HS$, in the context of MC, in comparison with that of the standard PTLs $\LTL$, $\CTL$, and $\CTLStar$. The analysis is carried on \emph{enforcing the homogeneity assumption}.
% 
We prove that $\HS$ endowed with the state-based semantics (hereafter denoted as $\HS_\stat$) is not comparable with $\LTL$, $\CTL$, and $\CTLStar$. On one hand, the result supports the intuition that $\HS_\stat$ gains some expressiveness by the ability of branching in the past. On the other hand, $\HS_\stat$ does not feature the possibility of forcing the verification of a property over an %linear 
infinite path, thus implying that the formalisms are not comparable. With the aim of having a more ``effective'' comparison base, we consider two additional semantic variants of $\HS$, 
%besides the state-based one $\HS_\stat$, 
namely, the \emph{computation-tree-based semantic variant} (denoted as $\HS_\LinearPast$) and the \emph{trace-based} one ($\HS_\LinearTime$). 
\input{Chaps/TOCL17/Imgs/semPicturesLN.tex}

The state-based (see Figure~\ref{fig:ST}) and computation-tree-based (see Figure~\ref{fig:CT}) approaches rely on a \emph{branching}-time setting and differ in the nature of past. In the latter approach, past is \emph{linear}: each interval may have several possible futures, but only a unique past. Moreover, past is assumed to be \emph{finite} 
%(since program computations have a definite starting time) 
and \emph{cumulative}, that is, the story of the current situation increases with time, and is never forgotten. 
%
%Finally, the 
The trace-based approach relies on a \emph{linear}-time setting (see Figure~\ref{fig:LN}), where the infinite paths (computations) of the given Kripke structure are the main semantic entities. Branching is neither allowed in the past nor in the future.
%
Note that the linear-past (rather than branching) approach is more suited to the specification  of dynamic behaviors, because it considers states in a computation tree, while the branching-past approach considers machine states, where past is not very meaningful for the specification of behavioral constraints~\cite{LS95}.

The variant $\HS_\LinearPast$ is a natural candidate for an expressiveness comparison with the branching time logics  $\CTL$ and $\CTLStar$. The most interesting and technically involved result is the characterization of the expressive power of $\HS_\LinearPast$: $\HS_\LinearPast$ turns out to be expressively equivalent to finitary $\CTLStar$, that is, the variant of $\CTLStar$ with quantification over \emph{finite} paths. As for $\CTL$, a non comparability result can be stated.
%
%Conversely, 
The variant $\HS_\LinearTime$ is a natural candidate for an expressiveness comparison with $\LTL$: 
%As a matter of fact, 
we prove that $\HS_\LinearTime$ and $\LTL$ are equivalent (this result holds true even for a very small fragment of $\HS_\LinearTime$), but the former is at least exponentially more succinct than the latter. 

We complete the picture with a comparison of the three semantic variants $\HS_\stat$, $\HS_\LinearPast$, and $\HS_\LinearTime$. We show that, as expected, $\HS_\LinearTime$ is not comparable with either of the branching versions, $\HS_\LinearPast$ and $\HS_\stat$. The interesting result is that, on the other hand, $\HS_\LinearPast$ is strictly included in $\HS_\stat$: this supports $\HS_\stat$ as a reasonable and adequate semantic choice.

\begin{figure}[b]
%\vspace*{-0.4cm}
\centering
%\resizebox{\width}{\height}{%
\begin{tikzpicture}[-,>=stealth',shorten >=1pt,auto,semithick,main node/.style={rectangle,draw,inner sep=2pt}]  
%
\tikzstyle{gray node}=[fill=gray!30]
%
\node [main node](0) at (0,0) {$\HS_\LinearPast$};
\node [main node](1) at (0,1.5) {$\HS_\LinearTime$};
\node [main node](2) at (0,-1.5) {$\HS_\stat$};
\node [main node](3) at (2.5,0) {finitary $\CTLStar$};
\node [main node](4) at (2.5,1.5) {$\LTL$};
\node [main node](6) at (2.5,-1.5) {$\CTL$};
\node [main node](5) at (5,0) {$\CTLStar$};
%   
\draw [dashed] (0.east) to node {$\equiv$} (3);
\draw [dashed] (1.east) to node {$\equiv$} (4);
\draw [dashed] (3.east) to node {$<$} (5);
\draw [dashed] (0.north) to node [swap] {$\neq$} (1);
\draw [dashed] (0.south) to node {\rotatebox{-90}{$<$}} (2);
\draw [dashed] (2.east) to node [swap] {$\neq$} (6.west);
\draw [dashed] (2) [out=-40,in=270] to node [near end] {$\neq$} (5);
\draw [dashed] (0.south east) to node {$\neq$} (6.west);
\draw [dashed] (2) [out=135,in=225] to node {$\neq$} (1);
\end{tikzpicture}%}
\vspace*{-0.6cm}
\caption{Overview of the expressiveness results.}\label{results}
\end{figure}

The complete picture of the expressiveness results is reported in Figure~\ref{results}
(the symbols $\neq$, $\equiv$, and $<$ denote incomparability, equivalence, and strict %expressiveness 
inclusion, respectively).

\paragraph*{Organization of the chapter.}
\begin{itemize}
	\item In the next section we start with some preliminaries; in particular, in Section~\ref{sect:PTL} we recall the well-known PTLs $\LTL$, $\CTL$ and $\CTLStar$; in Section~\ref{sect:3sem} we define the three semantic variants of $\HS$ ($\HS_\stat$, $\HS_\LinearPast$ and $\HS_\LinearTime$). In Section~\ref{subs:vendingMach} we provide a detailed example which gives an intuitive account of the three semantic variants and highlights their differences.
	\item In the next three sections we analyze and compare the expressiveness of these logics.  
In Section~\ref{sec:CharacterizezionOfLeniarTimeHS} we show the expressive equivalence of $\LTL$ and $\HS_\LinearTime$. Then, in Section~\ref{sec:characterizationHSLinearPast} we prove the equivalence of $\HS_\LinearPast$ and finitary $\CTLStar$. In Section~\ref{sect:allSems} we compare the expressiveness of 
$\HS_\stat$, $\HS_\LinearPast$ and $\HS_\LinearTime$.
\end{itemize}

\emptydoublepage
\include{relatedwork}
\emptydoublepage
%!TEX root = ./main.tex
%
% This file is part of the i10 thesis template developed and used by the
% Media Computing Group at RWTH Aachen University.
% The current version of this template can be obtained at
% <http://www.media.informatik.rwth-aachen.de/karrer.html>.

\chapter{Eigene Beitr\"age}
\label{ownwork} 

\mnote{Beginning: This is a margin note that---hopefully---spans multiple lines so I can see if it fits in well with the rest of the text.}Lorem ipsum dolor sit amet, consectetur adipisicing elit, sed do eiusmod tempor incididunt ut labore et dolore magna aliqua. Ut enim ad minim veniam, quis nostrud exercitation ullamco laboris nisi ut aliquip ex ea commodo consequat. Duis aute irure dolor in reprehenderit in voluptate velit esse cillum dolore eu fugiat nulla pariatur. Excepteur sint occaecat cupidatat non proident, sunt in culpa qui officia deserunt mollit anim id est laborum. Lorem ipsum dolor sit amet, consectetur adipisicing elit, sed do eiusmod tempor incididunt ut labore et dolore magna aliqua. Ut enim ad minim veniam, quis nostrud exercitation ullamco laboris nisi ut aliquip ex ea commodo consequat. Duis aute irure dolor in reprehenderit in voluptate velit esse cillum dolore eu fugiat nulla pariatur. Excepteur sint occaecat cupidatat non proident, sunt in culpa qui officia deserunt mollit anim id est laborum. Lorem ipsum dolor sit amet, consectetur adipisicing elit, sed do eiusmod tempor incididunt ut labore et dolore magna aliqua. Ut enim ad minim veniam, quis nostrud exercitation ullamco laboris nisi ut aliquip ex ea commodo consequat. Duis aute irure dolor in reprehenderit in voluptate velit esse cillum dolore eu fugiat nulla pariatur. Excepteur sint occaecat cupidatat non proident, sunt in culpa qui officia deserunt mollit anim id est laborum.

Lorem \mnote{After first word: This is a margin note that---hopefully---spans multiple lines so I can see if it fits in well with the rest of the text.}ipsum dolor sit amet, consectetur adipisicing elit, sed do eiusmod tempor incididunt ut labore et dolore magna aliqua. Ut enim ad minim veniam, quis nostrud exercitation ullamco laboris nisi ut aliquip ex ea commodo consequat. Duis aute irure dolor in reprehenderit in voluptate velit esse cillum dolore eu fugiat nulla pariatur. Excepteur sint occaecat cupidatat non proident, sunt in culpa qui officia deserunt mollit anim id est laborum.Lorem ipsum dolor sit amet, consectetur adipisicing elit, sed do eiusmod tempor incididunt ut labore et dolore magna aliqua. Ut enim ad minim veniam, quis nostrud exercitation ullamco laboris nisi ut aliquip ex ea commodo consequat. Duis aute irure dolor in reprehenderit in voluptate velit esse cillum dolore eu fugiat nulla pariatur. Excepteur sint occaecat cupidatat non proident, sunt in culpa qui officia deserunt mollit anim id est laborum.Lorem ipsum dolor sit amet, consectetur adipisicing elit, sed do eiusmod tempor incididunt ut labore et dolore magna aliqua. Ut enim ad minim veniam, quis nostrud exercitation ullamco laboris nisi ut aliquip ex ea commodo consequat. Duis aute irure dolor in reprehenderit in voluptate velit esse cillum dolore eu fugiat nulla pariatur. Excepteur sint occaecat cupidatat non proident, sunt in culpa qui officia deserunt mollit anim id est laborum.

Lorem ipsum dolor sit amet, consectetur adipisicing elit, sed do eiusmod tempor incididunt ut labore et dolore magna aliqua. Ut enim ad minim veniam, quis nostrud exercitation ullamco laboris nisi ut aliquip ex ea commodo consequat. Duis aute irure dolor in reprehenderit in voluptate velit esse cillum dolore eu fugiat nulla pariatur. Excepteur sint occaecat cupidatat non proident, sunt in culpa qui officia deserunt mollit anim id est laborum.
Lorem ipsum dolor sit amet, consectetur adipisicing elit, sed do eiusmod tempor incididunt ut labore et dolore magna aliqua. Ut enim ad minim veniam, quis nostrud exercitation ullamco laboris nisi ut aliquip ex ea commodo consequat. Duis aute irure dolor in reprehenderit in voluptate velit esse cillum dolore eu fugiat nulla pariatur. Excepteur sint occaecat cupidatat non proident, sunt in culpa qui officia deserunt mollit anim id est laborum. \mnote{End: This is a margin note that---hopefully---spans multiple lines so I can see if it fits in well with the rest of the text.}

Lorem ipsum dolor sit amet, consectetur adipisicing elit, sed do eiusmod tempor incididunt ut labore et dolore magna aliqua. Ut enim ad minim veniam, quis nostrud exercitation ullamco laboris nisi ut aliquip ex ea commodo consequat. \mnote{Middle: This is a margin note that---hopefully---spans multiple lines so I can see if it fits in well with the rest of the text.}Duis aute irure dolor in reprehenderit in voluptate velit esse cillum dolore eu fugiat nulla pariatur. Excepteur sint occaecat cupidatat non proident, sunt in culpa qui officia deserunt mollit anim id est laborum.

Lorem ipsum dolor sit amet, consectetur adipisicing elit, sed do eiusmod tempor incididunt ut labore et dolore magna aliqua. Ut enim ad minim veniam, quis nostrud exercitation ullamco laboris nisi ut aliquip ex ea commodo consequat. Duis aute irure dolor in reprehenderit in voluptate velit esse cillum dolore eu fugiat nulla pariatur. Excepteur sint occaecat cupidatat non proident, sunt in culpa qui officia deserunt mollit anim id est laborum.

Lorem ipsum dolor sit amet, consectetur adipisicing elit, sed do eiusmod tempor incididunt ut labore et dolore magna aliqua. Ut enim ad minim veniam, quis nostrud exercitation ullamco laboris nisi ut aliquip ex ea commodo consequat. Duis aute irure dolor in reprehenderit in voluptate velit esse cillum dolore eu fugiat nulla pariatur. Excepteur sint occaecat cupidatat non proident, sunt in culpa qui officia deserunt mollit anim id est laborum.

\begin{itemize}
	\item \mnote{in first bullet: This is a margin note that---hopefully---spans multiple lines so I can see if it fits in well with the rest of the text.}Lorem ipsum dolor sit amet, consectetur adipisicing elit,
	\item sed do eiusmod tempor incididunt ut labore et dolore magna aliqua
	\item Ut enim ad minim veniam, quis nostrud exercitation ullamco laboris nisi ut aliquip ex ea commodo consequat.
	\item[a] Excepteur sint occaecat cupidatat non proident, sunt in culpa qui.
	\item[b] Excepteur sint occaecat cupidatat non proident, sunt in culpa qua.
\end{itemize}

Lorem ipsum dolor sit amet, consectetur adipisicing elit, sed do eiusmod tempor incididunt ut labore et dolore magna aliqua. Ut enim ad minim veniam, quis nostrud exercitation ullamco laboris nisi ut aliquip ex ea commodo consequat. Duis aute irure dolor in reprehenderit in voluptate velit esse cillum dolore eu fugiat nulla pariatur. Excepteur sint occaecat cupidatat non proident, sunt in culpa qui officia deserunt mollit anim id est laborum.Lorem ipsum dolor sit amet, consectetur adipisicing elit, sed do eiusmod tempor incididunt ut labore et dolore magna aliqua. Ut enim ad minim veniam, quis nostrud exercitation ullamco laboris nisi ut aliquip ex ea commodo consequat. Duis aute irure dolor in reprehenderit in voluptate velit esse cillum dolore eu fugiat nulla pariatur. Excepteur sint occaecat cupidatat non proident, sunt in culpa qui officia deserunt mollit anim id est laborum.Lorem ipsum dolor sit amet, consectetur adipisicing elit, sed do eiusmod tempor incididunt ut labore et dolore magna aliqua. Ut enim ad minim veniam, quis nostrud exercitation ullamco laboris nisi ut aliquip ex ea commodo consequat. Duis aute irure dolor in reprehenderit in voluptate velit esse cillum dolore eu fugiat nulla pariatur. Excepteur sint occaecat cupidatat non proident, sunt in culpa qui officia deserunt mollit anim id est laborum.Lorem ipsum dolor sit amet, consectetur adipisicing elit, sed do eiusmod tempor incididunt ut labore et dolore magna aliqua. Ut enim ad minim veniam, quis nostrud exercitation ullamco laboris nisi ut aliquip ex ea commodo consequat. Duis aute irure dolor in reprehenderit in voluptate velit esse cillum dolore eu fugiat nulla pariatur. Excepteur sint occaecat cupidatat non proident, sunt in culpa qui officia deserunt mollit anim id est laborum.

Lorem ipsum dolor sit amet, consectetur adipisicing elit, sed do eiusmod tempor incididunt ut labore et dolore magna aliqua. Ut enim ad minim veniam, quis nostrud exercitation ullamco laboris nisi ut aliquip ex ea commodo consequat. Duis aute irure dolor in reprehenderit in voluptate velit esse cillum dolore eu fugiat nulla pariatur. Excepteur sint occaecat cupidatat non proident, sunt in culpa qui officia deserunt mollit anim id est laborum.

Lorem ipsum dolor sit amet, consectetur adipisicing elit, sed do eiusmod tempor incididunt ut labore et dolore magna aliqua. Ut enim ad minim veniam, quis nostrud exercitation ullamco laboris nisi ut aliquip ex ea commodo consequat. Duis aute irure dolor in reprehenderit in voluptate velit esse cillum dolore eu fugiat nulla pariatur. Excepteur sint occaecat cupidatat non proident, sunt in culpa qui officia deserunt mollit anim id est laborum.Lorem ipsum dolor sit amet, consectetur adipisicing elit, sed do eiusmod tempor incididunt ut labore et dolore magna aliqua. Ut enim ad minim veniam, quis nostrud exercitation ullamco laboris nisi ut aliquip ex ea commodo consequat. Duis aute irure dolor in reprehenderit in voluptate velit esse cillum dolore eu fugiat nulla pariatur. Excepteur sint occaecat cupidatat non proident, sunt in culpa qui officia deserunt mollit anim id est laborum.Lorem ipsum dolor sit amet, consectetur adipisicing elit, sed do eiusmod tempor incididunt ut labore et dolore magna aliqua. Ut enim ad minim veniam, quis nostrud exercitation ullamco laboris nisi ut aliquip ex ea commodo consequat. Duis aute irure dolor in reprehenderit in voluptate velit esse cillum dolore eu fugiat nulla pariatur. Excepteur sint occaecat cupidatat non proident, sunt in culpa qui officia deserunt mollit anim id est laborum.

Lorem ipsum dolor sit amet, consectetur adipisicing elit, sed do eiusmod tempor incididunt ut labore et dolore magna aliqua. Ut enim ad minim veniam, quis nostrud exercitation ullamco laboris nisi ut aliquip ex ea commodo consequat. Duis aute irure dolor in reprehenderit in voluptate velit esse cillum dolore eu fugiat nulla pariatur. Excepteur sint occaecat cupidatat non proident, sunt in culpa qui officia deserunt mollit anim id est laborum.

Lorem ipsum dolor sit amet, consectetur adipisicing elit, sed do eiusmod tempor incididunt ut labore et dolore magna aliqua. Ut enim ad minim veniam, quis nostrud exercitation ullamco laboris nisi ut aliquip ex ea commodo consequat. Duis aute irure dolor in reprehenderit in voluptate velit esse cillum dolore eu fugiat nulla pariatur. Excepteur sint occaecat cupidatat non proident, sunt in culpa qui officia deserunt mollit anim id est laborum.

\emptydoublepage
%!TEX root = ./main.tex
%
% This file is part of the i10 thesis template developed and used by the
% Media Computing Group at RWTH Aachen University.
% The current version of this template can be obtained at
% <http://www.media.informatik.rwth-aachen.de/karrer.html>.

\chapter{Evaluierung}
\label{evaluation}
\index{evaluierung|(}

Lorem ipsum dolor sit amet, consectetur adipisicing elit, sed do eiusmod tempor incididunt ut labore et dolore magna aliqua. Ut enim ad minim veniam, quis nostrud exercitation ullamco laboris nisi ut aliquip ex ea commodo consequat. Duis aute irure dolor in reprehenderit in voluptate velit esse cillum dolore eu fugiat nulla pariatur. Excepteur sint occaecat cupidatat non proident, sunt in culpa qui officia deserunt mollit anim id est laborum.Lorem ipsum dolor sit amet, consectetur adipisicing elit, sed do eiusmod tempor incididunt ut labore et dolore magna aliqua. Ut enim ad minim veniam, quis nostrud exercitation ullamco laboris nisi ut aliquip ex ea commodo consequat. Duis aute irure dolor in reprehenderit in voluptate velit esse cillum dolore eu fugiat nulla pariatur. Excepteur sint occaecat cupidatat non proident, sunt in culpa qui officia deserunt mollit anim id est laborum.Lorem ipsum dolor sit amet, consectetur adipisicing elit, sed do eiusmod tempor incididunt ut labore et dolore magna aliqua. Ut enim ad minim veniam, quis nostrud exercitation ullamco laboris nisi ut aliquip ex ea commodo consequat. Duis aute irure dolor in reprehenderit in voluptate velit esse cillum dolore eu fugiat nulla pariatur. Excepteur sint occaecat cupidatat non proident, sunt in culpa qui officia deserunt mollit anim id est laborum.Lorem ipsum dolor sit amet, consectetur adipisicing elit, sed do eiusmod tempor incididunt ut labore et dolore magna aliqua. Ut enim ad minim veniam, quis nostrud exercitation ullamco laboris nisi ut aliquip ex ea commodo consequat. Duis aute irure dolor in reprehenderit in voluptate velit esse cillum dolore eu fugiat nulla pariatur. Excepteur sint occaecat cupidatat non proident, sunt in culpa qui officia deserunt mollit anim id est laborum.

 \index{evaluierung|)}

\emptydoublepage
%!TEX root = ./main.tex
%
% This file is part of the i10 thesis template developed and used by the
% Media Computing Group at RWTH Aachen University.
% The current version of this template can be obtained at
% <http://www.media.informatik.rwth-aachen.de/karrer.html>.

\chapter{Zusammenfassung und Ausblick}
\label{summaryandfuturework}

%here
Lorem ipsum dolor sit amet, consectetur adipisicing elit, sed do eiusmod tempor incididunt ut labore et dolore magna aliqua. Ut enim ad minim veniam, quis nostrud exercitation ullamco laboris nisi ut aliquip ex ea commodo consequat. Duis aute irure dolor in reprehenderit in voluptate velit esse cillum dolore eu fugiat nulla pariatur. 

\section{Zusammenfassung und eigene Beitr\"age}
\label{summaryandfuturework.summary}

%here
Lorem ipsum dolor sit amet, consectetur adipisicing elit, sed do eiusmod tempor incididunt ut labore et dolore magna aliqua. Ut enim ad minim veniam, quis nostrud exercitation ullamco laboris nisi ut aliquip ex ea commodo consequat. Duis aute irure dolor in reprehenderit in voluptate velit esse cillum dolore eu fugiat nulla pariatur. Excepteur sint occaecat cupidatat non proident, sunt in culpa qui officia deserunt mollit anim id est laborum.Lorem ipsum dolor sit amet, consectetur adipisicing elit, sed do eiusmod tempor incididunt ut labore et dolore magna aliqua. Ut enim ad minim veniam, quis nostrud exercitation ullamco laboris nisi ut aliquip ex ea commodo consequat. Duis aute irure dolor in reprehenderit in voluptate velit esse cillum dolore eu fugiat nulla pariatur. Excepteur sint occaecat cupidatat non proident, sunt in culpa qui officia deserunt mollit anim id est laborum.Lorem ipsum dolor sit amet, consectetur adipisicing elit, sed do eiusmod tempor incididunt ut labore et dolore magna aliqua. Ut enim ad minim veniam, quis nostrud exercitation ullamco laboris nisi ut aliquip ex ea commodo consequat. Duis aute irure dolor in reprehenderit in voluptate velit esse cillum dolore eu fugiat nulla pariatur. Excepteur sint occaecat cupidatat non proident, sunt in culpa qui officia deserunt mollit anim id est laborum.Lorem ipsum dolor sit amet, consectetur adipisicing elit, sed do eiusmod tempor incididunt ut labore et dolore magna aliqua. Ut enim ad minim veniam, quis nostrud exercitation ullamco laboris nisi ut aliquip ex ea commodo consequat. Duis aute irure dolor in reprehenderit in voluptate velit esse cillum dolore eu fugiat nulla pariatur. Excepteur sint occaecat cupidatat non proident, sunt in culpa qui officia deserunt mollit anim id est laborum.

\section{Ausblick}
\label{summaryandfuturework.futurework}
\index{Ausblick|(}

%here
Lorem ipsum dolor sit amet, consectetur adipisicing elit, sed do eiusmod tempor incididunt ut labore et dolore magna aliqua. Ut enim ad minim veniam, quis nostrud exercitation ullamco laboris nisi ut aliquip ex ea commodo consequat. Duis aute irure dolor in reprehenderit in voluptate velit esse cillum dolore eu fugiat nulla pariatur. Excepteur sint occaecat cupidatat non proident, sunt in culpa qui officia deserunt mollit anim id est laborum.Lorem ipsum dolor sit amet, consectetur adipisicing elit, sed do eiusmod tempor incididunt ut labore et dolore magna aliqua. Ut enim ad minim veniam, quis nostrud exercitation ullamco laboris nisi ut aliquip ex ea commodo consequat. Duis aute irure dolor in reprehenderit in voluptate velit esse cillum dolore eu fugiat nulla pariatur. Excepteur sint occaecat cupidatat non proident, sunt in culpa qui officia deserunt mollit anim id est laborum.Lorem ipsum dolor sit amet, consectetur adipisicing elit, sed do eiusmod tempor incididunt ut labore et dolore magna aliqua. Ut enim ad minim veniam, quis nostrud exercitation ullamco laboris nisi ut aliquip ex ea commodo consequat. Duis aute irure dolor in reprehenderit in voluptate velit esse cillum dolore eu fugiat nulla pariatur. Excepteur sint occaecat cupidatat non proident, sunt in culpa qui officia deserunt mollit anim id est laborum.Lorem ipsum dolor sit amet, consectetur adipisicing elit, sed do eiusmod tempor incididunt ut labore et dolore magna aliqua. Ut enim ad minim veniam, quis nostrud exercitation ullamco laboris nisi ut aliquip ex ea commodo consequat. Duis aute irure dolor in reprehenderit in voluptate velit esse cillum dolore eu fugiat nulla pariatur. Excepteur sint occaecat cupidatat non proident, sunt in culpa qui officia deserunt mollit anim id est laborum.

\index{Ausblick|)}

\emptydoublepage

\begin{appendix}
\include{appendix_01}
\emptydoublepage
\include{appendix_02}
\emptydoublepage
\end{appendix}

%--------------------------------------------------------------
\backmatter

%Bibliographie
%bibliography
\clearpage
\phantomsection
\addcontentsline{toc}{chapter}{\protect\numberline{}Bibliography}
\bibliography{thesis}
\emptydoublepage

%Index
%index
\markboth{Index}{Index}
\thispagestyle{plain}
\phantomsection
\addcontentsline{toc}{chapter}{\protect\numberline{}Index}
\printindex

%Satzdatum
%date of typesetting
\newpage
\thispagestyle{empty}
\vspace*{\fill}
\hspace*{\fill}{\tiny Typeset \today}

%--------------------------------------------------------------
\end{document}

